\chapter{Introduction}
\section{Motivation}
Contrastive vision-language models map images and text into a shared embedding space, and from that single construction they inherit retrieval, zero-shot classification and a great deal else. The construction is also expensive. Model quality is coupled to large paired datasets, long training schedules and two towers optimised together, which places the method out of reach of most research settings and many deployment ones.

There is an obvious alternative. Strong vision encoders and strong sentence encoders already exist, trained separately and at someone else's expense; if the representations they have learned are compatible enough, aligning them should require only a small trainable module rather than a second round of web-scale pretraining. The expensive part has already happened. What remains is a bridge.

Whether that bridge can be built cheaply is not obvious, and this dissertation takes the question seriously enough to let it fail. Encoders trained independently need not expose compatible geometry at all. A training recipe that is dependable when both towers learn can behave quite differently when most parameters are held fixed, because the assumptions that justified it no longer hold. Efficiency claimed from parameter counts need not survive end-to-end measurement on real hardware. And model selection quietly breaks when baselines are imported from other checkpoints, other splits or other profiling sessions, which is common enough in this literature to be worth guarding against explicitly. Comparability and measurement are therefore treated here as part of the research problem rather than as reporting overhead.

\section{Research question}
\begin{quote}
\itshape
Can modern frozen self-supervised vision encoders, combined with compatibility-based pair selection and sub-five-million-parameter adaptation, approach compact jointly pretrained vision-language models under a fixed inference budget?
\end{quote}
The question is evaluated through a hierarchy rather than by relabelling one favourable baseline. MobileCLIP2-S0 is the primary compact jointly pretrained reference and also the distillation teacher. OpenCLIP~\cite{cherti2023scaling} ViT-B/32 is a widely used field anchor with a split-matched local validation score and same-allocation latency control; it is not described as compact. SigLIP2~\cite{tschannen2025siglip2} ViT-B/32 is an alternative-objective field reference. ``Approach'' means materially reducing a split-matched local accuracy gap while keeping locally trainable inference parameters below five million and satisfying a paired latency comparison. The dissertation reports a strictly frozen-backbone endpoint and a bounded LoRA~\cite{hu2021lora} adaptation endpoint. Both endpoints have since been evaluated once on the sealed Flickr30k~\cite{young2014flickr30k} test, so compact-reference retention is reported on matched splits rather than withheld as a cross-split comparison.

\section{Contributions}
The contributions are empirical and methodological rather than architectural.

\begin{enumerate}
\item A controlled reconstruction of the frozen-alignment problem that replaced mismatched historical baselines with local checkpoint- and protocol-matched references.
\item A diagnosis showing that a cross-batch memory queue, rather than the loss family, caused the dominant early performance collapse.
\item A batch, age, and modality factorial showing immediate stale-negative harm and demonstrating that drift magnitude alone does not determine the damage.
\item A constrained alignment frontier showing that learned selection over frozen patch/token sequences contributes substantially more than access to those sequences through mean pooling alone.
\item A bounded LoRA upper bound and a declared same-allocation timing amendment that separate model improvement from measurement-session variation.
\item A reproducible negative-result record: failed compatibility screening, harmful memory, weak web-data transfer under the tested budget, non-monotonic teacher quality, unsuccessful LR retuning, and hard-negative mining stopped by its frozen guard.
\item A closed speed--accuracy test showing that parameter-free internal token merging can improve measured latency while substantially damaging retrieval, so reduced token count is not by itself an efficiency-frontier improvement.
\end{enumerate}
\section{Thesis claim}
The evidence supports the following bounded answer. Modern frozen unimodal encoders can recover a substantial portion of the split-matched OpenCLIP field anchor's retrieval performance with fewer than five million locally trainable inference parameters and comparable measured latency, but compact-reference parity is not demonstrated. Learned token aggregation is the most effective fully frozen-backbone intervention. Limited encoder adaptation closes more of the field-anchor gap, while training recipe and measurement validity matter enough to reverse model-selection conclusions.

\chapter{Methodology}
\section{Research design}
The study uses staged decision gates. Each stage either locks a choice for the next stage, preserves a negative outcome, or stops a branch. This reduces the degree to which later observations silently change earlier success criteria. The canonical register assigns stable IDs to all 29 experiment families. All planned computation is closed: one branch is incomplete-archived, one is superseded, one is a valid stopped-by-gate result, and the remaining 26 are complete.

The main progression is:

\begin{enumerate}
\item reconstruct locally measured paired and frozen baselines;
\item screen encoder compatibility;
\item diagnose the failed recipe rather than relax its gate;
\item isolate queue capacity, age, batch, and modality;
\item test data coverage and transfer under matched budgets;
\item correct latency and FLOP accounting;
\item test distillation, resolution, teacher, and aggregation choices;
\item freeze the strongest fully frozen-backbone model;
\item run a bounded encoder-adaptation upper bound;
\item repair an invalid cross-session latency comparison in one declared same-allocation measurement;
\item profile fixed internal token merging; and
\item run the frozen three-seed accuracy closure and retain the no-merge models when every merge arm loses accuracy.
\end{enumerate}
\section{Models}
The final fully frozen-backbone model uses a DINOv3~\cite{simeoni2025dinov3} ViT-S/16 vision encoder at 224 pixels and all-MiniLM-L6-v2 as the text encoder. Vision patch tokens are aggregated by the C4 configuration: width 256 with two transformer blocks. Text tokens are pooled by a 128-dimensional, four-head learned query. Residual projections map both towers into a normalised 384-dimensional shared space.

\subsection{Formulation}
Write $f\_V$ and $f\_T$ for the vision and text encoders. Both are held at their pretrained values, so for an image $x$ and a caption $y$ the token sequences

\begin{equation*}
H_V = f_V(x) \in \mathbb{R}^{N_V \times d_V},
\qquad
H_T = f_T(y) \in \mathbb{R}^{N_T \times d_T}
\end{equation*}
are fixed functions of the input: no gradient reaches $f\_V$ or $f\_T$. Everything the alignment can learn therefore has to live in what reads those sequences.

That reader is a learned-query attention block. Given $k$ learned queries $Q \in \mathbb{R}^{k \times d}$ and a frozen token sequence $H$, the aggregator computes

\begin{equation*}
A(H) = \mathrm{softmax}\!\left(
  \frac{(QW_q)(HW_k)^{\top}}{\sqrt{d_h}}
\right) HW_v ,
\end{equation*}
with $W\_q, W\_k, W\_v$ trainable. The queries, not the tokens, decide what is read out. Mean pooling is the special case in which the attention weights are held uniform and nothing is selected, which is why the comparison between the two in Section 4.6 isolates selection rather than access.

Each aggregated representation is mapped into the shared space by a residual projection $\pi$ and $\ell\_2$-normalised,

\begin{equation*}
z_V = \frac{\pi_V(A_V(H_V))}{\lVert \pi_V(A_V(H_V)) \rVert_2},
\qquad
z_T = \frac{\pi_T(A_T(H_T))}{\lVert \pi_T(A_T(H_T)) \rVert_2},
\qquad
z_V, z_T \in \mathbb{R}^{384},
\end{equation*}
so that the similarity of an image and a caption is the inner product $s\_{ij} = z\_{V,i}^{\top} z\_{T,j}$.

Training minimises a symmetric InfoNCE objective over a batch $B$ with learned temperature $\tau$. Because COCO~\cite{lin2014coco} images carry several captions, the positive set $P(i)$ for image $i$ contains every in-batch caption sharing its image identifier rather than a single diagonal entry:

\begin{equation*}
\mathcal{L}_{V \rightarrow T}
= -\frac{1}{|B|} \sum_{i \in B} \log
\frac{\sum_{p \in P(i)} \exp(s_{ip} / \tau)}
     {\sum_{j \in B} \exp(s_{ij} / \tau)},
\qquad
\mathcal{L} = \tfrac{1}{2}\left(
  \mathcal{L}_{V \rightarrow T} + \mathcal{L}_{T \rightarrow V}
\right).
\end{equation*}
A cross-batch memory queue changes exactly one thing in this expression: it extends the denominator with embeddings $\tilde{z}$ that were produced by *earlier* versions of the projector and stored,

\begin{equation*}
\sum_{j \in B} \exp(s_{ij} / \tau)
\;\longrightarrow\;
\sum_{j \in B} \exp(s_{ij} / \tau)
+ \sum_{m \in M} \exp\!\left(z_{V,i}^{\top} \tilde{z}_{T,m} / \tau\right),
\end{equation*}
where $M$ is the stored memory. The encoders being frozen does not make $\tilde{z}$ safe, because $\tilde{z}$ is a function of the projector, and the projector is still moving. Chapter 4 is largely an account of what that substitution costs.

FreezeShift relaxes the frozen constraint by the smallest amount that could still be called an adaptation. For a base weight matrix $W$ it learns a rank-$r$ update

\begin{equation*}
W' = W + \frac{\alpha}{r} BA,
\qquad
B \in \mathbb{R}^{d \times r},
\quad
A \in \mathbb{R}^{r \times k},
\quad
r = 128,
\end{equation*}
with $W$ fixed. Since $W'$ is formed once and merged before deployment, the adapted model has the same inference graph as the frozen one, which is what makes the latency comparison in Section 4.8 a like-for-like measurement.

MobileCLIP2-S0 supplies a training-only distillation signal. It is absent at inference. Training uses queue-free InfoNCE, all available COCO captions, batch 1024, and a 24-epoch cosine schedule. \MTone{} contains 2,896,389 trainable inference parameters in a 47,196,549-parameter deployed stack; the underlying encoder parameters remain frozen.

FreezeShift adds rank-128 LoRA to the query/key/value/output attention projections in the final four DINOv3 blocks, the query/value projections in the final four MiniLM blocks, or both. The base weights remain fixed and LoRA is merged before latency profiling. Vision-, text-, and dual-tower totals are 4,076,037, 3,682,821, and 4,862,469 trainable inference parameters. The corresponding full deployed parameter counts are 48,376,197, 47,982,981, and 49,162,629. ``Trainable'' and ``total'' are reported separately throughout.

\section{Data and split policy}
COCO training data supplies paired alignment examples. COCO val2017 is treated as development evidence because it was involved in earlier model selection. Flickr30k Karpathy validation supplies the primary external development metric for the final alignment line. Its test split was held sealed throughout model development and opened once, for the confirmatory evaluation reported later in this dissertation. The 224px distilled configuration and checkpoint epochs were both chosen with exposure to Flickr validation. The measured difference between COCO-dev-selected and Flickr-selected epochs is 0.128pp; this quantifies one part of, but does not remove, the optimistic selection exposure.

Classification development analyses use deterministic partitions of CIFAR-100 train, Oxford-IIIT Pets trainval, and a stored stratified EuroSAT split. Official classification tests are reserved. Historical evaluations that touched official tests are retained as exploratory evidence and excluded from confirmatory claims. Split roles and leakage rules are machine-readable in \texttt{configs/\linebreak[1]evaluation\_\linebreak[1]protocol.\linebreak[1]yaml}.

\section{Metrics and selection}
The primary accuracy metric is mean bidirectional Recall@1:

\begin{equation*}
R@1_{mean}=\frac{R@1_{image\rightarrow text}+R@1_{text\rightarrow image}}{2}.
\end{equation*}
Final candidate scores are three-seed means, with standard deviations reported. FreezeShift promotion required the best parameter- and latency-eligible arm to improve over \MTone{} by more than one pooled standard deviation. This is a practical heuristic, not a null-hypothesis significance test.

Latency is full-stack Q3 time at batch 64 in native bfloat16 on an NVIDIA RTX PRO 6000 Blackwell. The corrective experiment used ten independently warmed repetitions, 100 timed iterations per repetition, and deterministic randomised order within each repetition. Candidate-minus-OpenCLIP paired differences are the primary comparison; bootstrap intervals classify faster, equivalent, or slower behaviour.

\section{Queue factorial}
The queue programme separates no-queue, capacity, nominal age, realised eviction age, batch size, and image-only versus text-only memory. Representation drift at eviction is measured rather than inferred from the nominal setting.

Drift is instrumented on a fixed probe of 512 held-out development pairs whose frozen encoder features are cached once, so that any movement observed is movement of the projector and nothing else. Every ten optimizer steps the probe is re-projected and normalised, and the step distance

\begin{equation*}
\delta(t) = \frac{1}{|S|} \sum_{s \in S}
\left( 1 - z_s^{(t)\top} z_s^{(t - \Delta)} \right),
\qquad \Delta = 10,
\end{equation*}
is recorded for each modality. The quantity reported at eviction is not the displacement between two endpoints but the *path length* accumulated while an entry actually sat in the queue: for a measured residence of $R$ optimizer steps,

\begin{equation*}
D(t; R) = \sum_{\text{intervals } \subseteq (t - R,\, t]} \delta(\cdot),
\end{equation*}
with boundary intervals prorated by overlap. Path length is the conservative choice, because a representation can wander a long way and return, and an endpoint measure would score that as no drift at all.

Damage is expressed against the matched queue-free control trained under otherwise identical settings,

\begin{equation*}
\mathrm{degradation} = 100 \times
\left( \mathrm{R@1}^{\text{no queue}}_{\text{mean}}
     - \mathrm{R@1}^{\text{queue}}_{\text{mean}} \right)
\ \text{pp},
\end{equation*}
so every degradation figure in Chapter 4 is a within-condition contrast rather than a comparison against a distant baseline. An age-zero row identifies the residual batch/in-batch-negative-count change. Predictions and gates were recorded before outcome analysis; the preregistered modality-asymmetry prediction was scored as a miss when its nominal-axis form failed.

\section{Reproducibility and integrity}
Configs, compact reports, fingerprints, manifests, and integrity receipts are retained in the submission. Heavyweight checkpoints, caches, datasets, and raw Slurm logs remain recoverable in a locally manifested quarantine rather than in GitHub. Negative and stopped branches remain in the register. The environment is represented by direct dependency specifications and a verified version lock. The codebase includes unit, integration, schema, protocol, and artifact tests.

\chapter{Results}
This chapter reports the experiments in the order the decision gates were passed, because several later choices only make sense once an earlier measurement has been corrected. The sequence has a shape. It opens with a failure that looked like a ceiling on frozen alignment and turned out to be a fault in the recipe (Sections 4.1 to 4.2); it then pursues the component responsible far enough to say something general about it (Section 4.3); and it spends the remainder establishing how much performance the repaired recipe can actually reach under the stated budget, through data, aggregation, bounded adaptation and token reduction (Sections 4.4 to 4.11).

Three of those later results are negative, and they are reported at the same length as the positive ones. A branch that was stopped by its own preregistered gate, a data source that failed to converge inside its compute ceiling, and a latency optimisation that worked and was rejected anyway are all part of the evidence rather than omissions from it.

\section{Foundations and the completed adaptive branch}
Early Phase 1 experiments screened five vision and four text encoders across baseline and bidirectional latent fusion variants, 23 configurations in total. The strongest grid result was DINOv2~\cite{oquab2023dinov2} ViT-S/14 with MiniLM under the local BLF variant, at 17.78\% mean COCO R@1 (19.52\% image-to-text, 16.05\% text-to-image), against 6.39\% for the weakest pair. Matched-seed BLF gains were not statistically stable: across the three matched baseline-versus-BLF comparisons, fusion won 7 of 15 metric contests and lost 8.

The more consequential result was that no pairing won everywhere. Scored per task the leaderboard reshuffles completely, and even the best *text* encoder changes with the task:

\begin{table}[htbp]
\centering
\caption{Best local encoder pairing by development task in the compatibility screen.}
\label{tab:compatibility}
\begin{tabular}{L{4.0cm}L{5.0cm}r}
\toprule
Task & Best pairing & Score \\
\midrule
COCO image-to-text R@1 & DINOv2 + MiniLM (local) & 19.52\% \\
COCO text-to-image R@1 & DINOv2 + MiniLM (local) & 16.05\% \\
CIFAR-100 zero-shot & ConvNeXtV2-T + MiniLM & 29.51\% \\
EuroSAT zero-shot & DINOv2 + BGE & 22.29\% \\
Oxford-IIIT Pets zero-shot & DINOv2 + BGE & 6.87\% \\
\bottomrule
\end{tabular}
\end{table}
Phase 1.5 then took the same 23 configurations to sample level, measuring cross-task rankings, pairwise complementarity over 759 configuration pairs and 37,191 per-class rows, a deliberately non-deployable best-of-all-models oracle across 6,075 subsets, and per-encoder efficiency profiles. Because no single pairing dominated, four diverse paths were carried forward rather than one. Two planned compositional evaluations, Winoground and SugarCrepe, were never run: the datasets were not installed, and they are recorded as missing rather than dropped, which is why compositional ability remains untested in Section 6.2. The Pets figure is also worth noting in hindsight, because the weakness it shows at 6.87\% is the same one the sealed evaluation finds in the finished models six weeks later.

Phase 2 trained cross-attention rerankers and a task-conditioned router. The branch completed 1,229 long-format result rows, 40 seed rows, nine statistical summaries, and four Pareto entries. Cross-attention improved some classification-development results while reducing COCO retrieval. Smaller candidate sets performed best, random negatives unexpectedly beat hard negatives, and the dense router's mean sample utility was 0.01436 below the best bridge with zero wins across five seeds. This mixed outcome is retained as a completed secondary branch, not folded into the final dual-encoder headline.

\section{Recipe diagnosis}
The initial compatibility screen did not clear its frozen gate. The natural reading of that failure was the discouraging one: that independently pretrained encoders simply do not expose compatible enough geometry, and that the ceiling had been found. The natural response would have been to lower the gate.

Neither was correct. Rather than move the threshold, Wave 0 decomposed the recipe and tested its components separately, and the result relocated the problem entirely. Removing the memory queue accounted for roughly 19 percentage points of improvement, whereas loss-family differences were worth about one. The ceiling had not been a property of the encoders at all. It had been a property of the training recipe, and specifically of a component included to help. Queue-free InfoNCE with all captions, batch 1024, and LR scale 3 was therefore locked.

Under this corrected recipe, the text-encoder ranking inverted: E5 exceeded BGE, which exceeded MiniLM. DINOv3 ViT-S/16 remained the preferred vision encoder. MiniLM was retained for continuity and efficiency; the ranking inversion itself is a result and constrains claims that encoder compatibility is intrinsic and recipe-independent.

\section{Cross-batch memory damage}
The capacity sweep showed monotonic harm:

\begin{table}[htbp]
\centering
\caption{Retrieval degradation as stale cross-batch memory capacity increases.}
\label{tab:queue}
\begin{tabular}{L{7.5cm}r}
\toprule
Queue condition & Mean R@1 \\
\midrule
No queue & 35.76\% \\
Increasing stale-memory capacity & monotonically decreasing \\
Capacity 16,384 & 16.77\% \\
\bottomrule
\end{tabular}
\end{table}
The queue supplied no favourable accuracy, memory, or throughput trade in this regime. The factorial located harm at measured age one: relative to matched queue-free controls, the loss was 6.46pp at batch 1024 and 5.60pp at batch 512. There was no observed safe non-zero age.

Drift was not a complete explanation. Four matched-drift comparisons remained within the declared tolerance while their degradation differed by as much as 4.74pp. Separate descriptive regressions produced an image slope of 5.553 and a text slope of 131.442 percentage points per unit drift. Because modality ranges differed, the conservative comparison restricted both to their overlapping drift interval; the resulting slopes were 21.4 and 153.3, a 7.2-fold ratio.

This does not identify a mechanism, and the dissertation does not claim one. False-negative structure, local geometry, and gradient sensitivity all remain possible contributors, and separating them would require instrumentation this study does not have. What the factorial does establish is a boundary that matters for practice: small cosine movement cannot by itself certify that a stored cross-modal negative is harmless, because equal movement demonstrably has modality-dependent consequences. The slow-drift argument is not refuted in general. It is shown to be insufficient here, which is enough to make any queue justified by drift alone an unsafe default in this regime.

One false-negative subtype can be excluded from the existing instrumentation. The loss builds its multi-positive mask after queued entries are appended, so a queued caption/image sharing the current \texttt{image\_\linebreak[1]id} is treated as positive, not negative. Non-zero queue-positive fractions at epoch boundaries confirm that this path was exercised. Unannotated semantic matches across different image IDs remain possible and are the residual false-negative hypothesis.

\section{Data coverage and transfer}
At matched update count, full COCO exceeded a 25\% subset by 15.03pp on Flickr validation, showing that unique in-distribution coverage mattered. A qualified CC3M mirror then underperformed the COCO control by 8.13pp at matched compute and by 5.89pp after four passes. The extended arm reached its compute ceiling without convergence, so this is not a universal claim about CC3M.

Mixing COCO back into CC3M training recovered much of the transfer loss: 53.133\% Flickr-validation mean R@1 versus 47.209\% for CC3M-only. Under this budget, source and caption distribution mattered more than raw pair count.

\section{Measurement corrections, resolution, and distillation}
The first efficiency frontier exposed three conclusion-changing artifacts. COCO could not be the primary final metric because selection had touched it; dynamic rather than fixed text padding was required for deployed timing; and teacher fusion and operator-level FLOP accounting had to match the actual inference graph.

Distillation had in fact been scoped before this point. A bounded capture probe had run MobileCLIP2-S0 and SigLIP2-B/32 as separate matched teacher branches over two seeds, with a preregistered rule that read the absolute student score rather than the capture fraction, and it returned a limited but measurable signal. That outcome is why the programme pursued efficiency-normalised evidence rather than chasing absolute retrieval through distillation alone.

At 224 pixels, DINOv3 closed the latency gap without learned-position interpolation mismatch. A 24-epoch trajectory quantified the difference between COCO-dev and Flickr-validation epoch selection. MobileCLIP2~\cite{faghri2025mobileclip2} distillation continued to help after the recipe was corrected, showing that its benefit was not merely compensation for the queue. Teacher benchmark strength was not monotonic with student transfer, and an equal two-teacher mixture did not beat MobileCLIP2 alone. Tuned mixtures were not tested.

\section{Token aggregation}
The controlled pooling comparison was decisive:

\begin{table}[htbp]
\centering
\caption{Effect of vision-token aggregation strategy, measured against a matched baseline.}
\label{tab:aggregation}
\begin{tabular}{L{7.5cm}r}
\toprule
Vision-token intervention & Gain over matched baseline \\
\midrule
Mean pooling & about +1.33pp \\
Learned-query aggregation & about +6.08pp \\
Transformer aggregation & about +6.91pp \\
\bottomrule
\end{tabular}
\end{table}
Thus patch access alone explained little of the gain. A learned selection or interaction mechanism was required. The 2$\times$2 width/depth study selected C4 (width 256, two blocks). Width and depth both helped but composed sub-additively, and further capacity was not justified below the five-million budget.

Text-side learned-query aggregation produced \MTone{} at 54.487\% mean R@1. Final LR factor 1.5 lost 0.763pp and failed its promotion rule, retaining factor 1.0. Guarded hard-negative mining was stopped before training because its frozen safety rule declared the candidate set infeasible.

\section{Fully frozen-backbone endpoint}
\MTone{} is the strongest fully frozen-backbone result:

\begin{table}[htbp]
\centering
\caption{Properties of \MTone{}, the strongest fully frozen-backbone model.}
\label{tab:mt1}
\begin{tabular}{L{8.5cm}r}
\toprule
Property & \MTone{} \\
\midrule
Flickr30k-validation mean bidirectional R@1 & 54.487\% \\
Three-seed SD & 0.328pp \\
Trainable inference parameters & 2,896,389 \\
Same-allocation mean Q3 & 9.148 ms \\
Paired Q3 difference vs OpenCLIP & -0.119 ms \\
95\% bootstrap interval & [-0.139, -0.080] ms \\
\bottomrule
\end{tabular}
\end{table}
It is therefore faster than the local OpenCLIP reference under the corrected same-allocation protocol. Its accuracy remains 15.454pp lower on validation.

\section{FreezeShift upper bound}
Every LoRA arm materially improved validation accuracy:

\begin{table}[htbp]
\centering
\caption{FreezeShift low-rank adaptation arms against the frozen \MTone{} baseline.}
\label{tab:freezeshift}
\begin{tabular}{L{4.5cm}rrr}
\toprule
Arm & Validation R@1 & Gain vs \MTone{} & Parameters \\
\midrule
Vision & 58.047\% & +3.560pp & 4.076M \\
Text & 59.724\% & +5.237pp & 3.683M \\
Dual & \textbf{63.416\%} & \textbf{+8.928pp} & \textbf{4.862M} \\
\bottomrule
\end{tabular}
\end{table}
The first reporter rejected all arms against a 9.480-ms ceiling taken from a different session. In that allocation the unchanged \MTone{} control itself measured 9.497 ms and failed, demonstrating that the absolute gate was contaminated by session variation. The negative verdict was preserved and a measurement amendment was declared before corrective profiling.

In the corrective allocation, OpenCLIP measured 9.267 ms. \MTone{}, vision, text, and dual measured 9.148, 9.168, 9.173, and 9.166 ms respectively. Every paired candidate-minus-reference interval lay below zero. The dual arm therefore satisfies the unchanged parameter and improvement rules under the repaired comparison and is the final development adaptation candidate.

Dual retains 90.67\% of OpenCLIP validation R@1 and closes 57.77\% of the \MTone{}-to-OpenCLIP gap. It does not beat OpenCLIP and is not a fully frozen encoder result.

\section{TokenShift profile}
TokenShift reduced 196 patch tokens to 49 after either block 8 or block 6 while preserving CLS and four register tokens:

\begin{table}[htbp]
\centering
\caption{Measured latency effect of fixed internal token merging.}
\label{tab:tokenshift-latency}
\begin{tabular}{L{5.5cm}rr}
\toprule
Arm & Mean Q3 & Paired change vs no merge \\
\midrule
No merge & 9.122 ms & --- \\
Merge after block 8 & 8.138 ms & -0.984 ms \\
Merge after block 6 & 7.515 ms & -1.606 ms \\
\bottomrule
\end{tabular}
\end{table}
Both paired intervals were wholly below zero. Because no accuracy evaluation was part of this profiling stage, the result established only that internal token count was latency-relevant and motivated the frozen accuracy closure.

\section{TokenShift accuracy closure}
The follow-up crossed the two merge depths with the strictly frozen \MTone{} and dual-LoRA parents, using three seeds per arm. Every merge remained faster, but none preserved its matched parent's retrieval accuracy:

\begin{table}[htbp]
\centering
\caption{Accuracy and latency of token merging against matched no-merge parents.}
\label{tab:tokenshift-accuracy}
\begin{tabular}{L{4.5cm}rrr}
\toprule
Parent and merge & Validation R@1 & Change vs parent & Mean Q3 \\
\midrule
Frozen, block 8 & 41.907\% & -12.581pp & 8.204 ms \\
Frozen, block 6 & 38.537\% & -15.950pp & 7.573 ms \\
Dual LoRA, block 8 & 57.725\% & -5.690pp & 8.203 ms \\
Dual LoRA, block 6 & 52.554\% & -10.861pp & 7.556 ms \\
\bottomrule
\end{tabular}
\end{table}
The later merge was consistently less damaging, and LoRA recovered part of the loss, but even the best trade lost 5.690pp for about 0.934 ms. Because the package prohibited automatic promotion, the final decision was recorded after all arms closed: retain frozen \MTone{} and dual FreezeShift without token merging. The result converts the earlier profile into a negative frontier finding---fewer internal tokens were faster, but not more efficient once accuracy was included.

\section{Reference and peer comparison}
Reference results must be read in protocol blocks. The final-candidate rows below use Flickr30k validation; the compact and SigLIP2 rows retained by the historical frontier use Flickr30k test. They are shown together to document coverage, not to create a cross-split ranking.

\begin{table}[htbp]
\centering
\caption{Final candidates against the compact reference, field anchor and alternative-objective reference. Protocol blocks differ: candidate rows use Flickr30k validation; the compact and SigLIP2 rows use Flickr30k test.}
\label{tab:references}
\small
\begin{tabular}{L{2.1cm}L{2.9cm}R{1.9cm}R{2.3cm}L{3.3cm}}
\toprule
Model & Role & Total parameters & Locally trainable parameters & Flickr evidence \\
\midrule
\MTone{} & strictly frozen endpoint & 47.197M & 2.896M & 54.487\% validation mean R@1 \\
Dual FreezeShift & bounded PEFT endpoint & 49.163M & 4.862M & 63.416\% validation mean R@1 \\
MobileCLIP2-S0 & primary compact reference and teacher & 74.835M & not applicable & 78.240\% test mean R@1 \\
OpenCLIP ViT-B/32 & field anchor & 151.277M & not applicable & 69.941\% validation; 68.220\% test \\
SigLIP2 ViT-B/32 & alternative-objective field reference & 376.856M & not applicable & 80.580\% test mean R@1 \\
\bottomrule
\end{tabular}
\end{table}
Dual FreezeShift is therefore 32.5\% of OpenCLIP's total size and reaches 90.67\% of its split-matched validation score. Against MobileCLIP2-S0 it is 65.7\% of the total size, not one third. A performance-retention ratio is not taken from this table, because doing so would divide validation by test; the retention figures against both references are reported from the sealed test results instead.

The direct frozen-alignment peer group is also necessary. Freeze-Align~\cite{maniparambil2025freezealign} reports 87.5/74.1 I2T/T2I R@1 for its large final configuration; the SAIL~\cite{zhang2025sail} supplement's common DINOv2-B/CC12M comparison reports 74.2/60.6 for SAIL-B and 68.1/49.3 for ShareLock~\cite{ruthardt2026sharelock}. STRUCTURE~\cite{groger2025structure} evaluates geometry-preserving alignment under limited data but no directly commensurate final Flickr configuration was identified. These published values use different encoders, data volumes, objectives and evaluation choices. They show that the present scores are not state of the art; they do not invalidate the queue mechanism study or the sub-five-million local-adaptation frontier.

\section{Where the residual gap lies}
Knowing that a gap to OpenCLIP remained said nothing about its character. A frozen-gallery ranking probe decomposed it, running the three C4 seeds and a matched OpenCLIP control over the identical Flickr30k validation gallery under one tie-break rule, with a determinism receipt confirming that repeated runs reproduced the same ranks. It is a diagnostic: the sealed test was not touched.

The decomposition splits every query into three outcomes---correct at rank 1, correct but ranked second to tenth, and not retrieved in the top ten at all. The distinction matters because the second class is a precision failure on material the model has already found, whereas the third is a failure to find it.

Both models fail in the same way, and the alignment head simply fails more often. Image-to-text, 29.8\% of queries land in ranks 2--10 against OpenCLIP's 18.2\%, while only 9.0\% fall outside the top ten against OpenCLIP's 1.7\%. Text-to-image, the split is 36.9\% against 30.1\% and 16.7\% against 10.1\%. Roughly three quarters of the image-to-text shortfall is therefore mis-ordering rather than missing evidence.

This locates the deficit without explaining it. The probe records that Flickr's five-caption structure inflates image-to-text relative to text-to-image for any dual encoder, so only the excess over the matched control carries a model-specific reading, and its ambiguity bins estimate an ambiguity-associated ceiling rather than an absolute one. What survives those caveats is a direction for future work: the remaining gap is mostly a ranking problem among plausible candidates, which is the regime where reranking and harder negatives act, rather than a representational failure to retrieve.

\section{Sealed final evaluation}
The final evaluation was run only after the student roster, selected epochs, checkpoint paths, prompts and dataset roles had been frozen. Its package exposed no optimiser or parameter-update path. Flickr30k Karpathy test was the primary confirmatory retrieval evaluation; CIFAR-100, Oxford-IIIT Pets and EuroSAT were supplementary zero-shot diagnostics with prior project exposure and are not presented as pristine confirmatory tests.

\begin{table}[htbp]
\centering
\caption{Frozen-roster Flickr30k test results for the selected endpoints and fixed references.}
\label{tab:final-test}
\begin{tabular}{L{4.5cm}rrr}
\toprule
Model & Flickr30k test mean R@1 & Seeds & Full-stack parameters \\
\midrule
\MTone{}, strictly frozen & 52.90 $\pm$ 0.87\% & 3 & 47.197M \\
Dual FreezeShift & 62.20 $\pm$ 0.45\% & 3 & 49.163M \\
OpenCLIP ViT-B/32 & 68.22\% & fixed reference & 151.277M \\
MobileCLIP2-S0 & 78.25\% & fixed reference & 74.835M \\
SigLIP2 ViT-B/32 & 80.46\% & fixed reference & 376.856M \\
\bottomrule
\end{tabular}
\end{table}
The 9.30pp difference between the two student endpoints is the directly measured test-set cost of strict freezing in this system. Validation-to-test declines were modest---1.59pp for \MTone{} and 1.22pp for FreezeShift---so the principal model ordering survived the sealed evaluation. Dual FreezeShift retained 91.18\% of the OpenCLIP field anchor with 32.5\% of its total parameters, but only 79.49\% of MobileCLIP2-S0 with 65.7\% of its parameters. It therefore approached the field anchor but remained well short of the compact jointly pretrained reference.

The classification diagnostics reinforced that boundary. The students were weak relative to all three references, especially on Oxford-IIIT Pets, and LoRA did not improve every diagnostic. COCO-trained retrieval alignment therefore did not reproduce broad CLIP-like semantic generalisation.

\chapter{Discussion}
Two kinds of result came out of the programme, and they deserve different weight. One is an engineering answer: how far frozen encoders can be pushed under a five-million-parameter budget, which is useful but bounded to this setting. The other is a finding about a widely used training component that was never specific to frozen alignment at all, and which the frozen setting merely made visible. This chapter argues that the second is the more durable contribution, sets out what neither result licenses, and treats the measurement corrections as findings rather than as housekeeping.

\section{Answer to the research question}
The answer is split by reference tier. Against the split-matched OpenCLIP field anchor, ``yes, in a bounded sense'': \MTone{} preserved 77.90\% of the validation score while meeting the parameter and latency constraints, and bounded dual-tower LoRA preserved 90.67\% with 4.862 million locally trainable inference parameters while measuring faster in the same allocation. The sealed Flickr30k test has since been opened, and it confirms the ordering: \MTone{} reaches 52.90 $\pm$ 0.87\% and dual FreezeShift 62.20 $\pm$ 0.45\% mean bidirectional R@1 across three seeds, retaining 91.18\% of the OpenCLIP field anchor but only 79.49\% of MobileCLIP2-S0. Against the actual compact reference, therefore, an accuracy ``approach'' claim is still not demonstrated: neither endpoint establishes parity, and the 9.30-point difference between them is the measured cost of strict freezing.

The progression also shows that architecture alone is an incomplete account. Queue removal produced a larger change than early module choices; data source and coverage mattered; teacher strength did not order student performance; and an invalid timing comparison initially reversed eligibility. In constrained systems, the recipe and measurement protocol are part of the model.

\section{Why the queue result matters}
XBM~\cite{wang2020xbm}'s slow-drift argument is plausible in the regime it studied, and nothing here contradicts it there. The claim under test is narrower and more practical: that slow movement is sufficient grounds for treating a stored negative as usable. It is not. The boundary appears once historical negatives cross independently parameterised modalities whose projectors continue to move. The observation that harm starts at age one makes ``use only a small queue'' unsupported here. The matched-drift dissociation further shows why monitoring one scalar movement threshold is insufficient.

The strongest objection is the known modality gap. Different cones may make a unit cosine change more consequential in one modality. But that does not erase the operational finding; it explains why drift-only safety rules fail. The descriptive data are also slope-dominated rather than a simple constant offset. The mechanism remains open, and the dissertation does not claim otherwise.

\section{What token aggregation demonstrates}
The token result should be read as evidence about frozen representations, not as architecture novelty. Global summaries hide useful information. Simply averaging patches recovers little, while learned interaction recovers much more. This is consistent with attentive-probing literature: frozen transformer features distribute task-relevant information across tokens, and a small learned reader can expose it efficiently.

\section{Adaptation versus freezing}
FreezeShift narrows the reference gap substantially. Its tower ablation is also informative: text-only adaptation gains more than vision-only, and dual adaptation gains most, so the residual mismatch is not isolated to one tower. At the same time, allowing LoRA changes the scientific category. The result is best interpreted as an upper bound on recoverable alignment under the same parameter budget, while \MTone{} remains the answer to the stricter frozen-backbone question.

This result must also be placed behind HALoRA~\cite{zhang2026halora} and related adapter/LoRA work in the novelty hierarchy. The useful evidence is the controlled vision/text/dual tower comparison under one local budget, not the invention of low-rank vision--language adaptation.

\section{Measurement as a research result}
The latency amendment is not cosmetic. A gate that rejects its unchanged control cannot support a model comparison. Preserving the original verdict, declaring the defect after observation, freezing the corrective protocol before measurement, and reporting paired uncertainty provides a transparent repair. It is weaker than a perfectly preregistered first attempt but stronger than silently replacing the threshold.

\section{Why TokenShift did not move the frontier}
TokenShift tests the final inference-side implication of the aggregation work. Earlier experiments showed that patch detail is valuable only when a learned reader can select it; TokenShift then removed three quarters of the patch tokens partway through the frozen vision tower. Its speed gain confirms that those tokens carry real computational cost, while its accuracy loss confirms that they still carry information used by the downstream reader. The milder loss after block 8 suggests later representations are more compressible, and the smaller loss under dual LoRA suggests adaptation can absorb some compression, but neither observation supports deployment of the tested fixed merge.

\chapter{Limitations and threats to validity}
\section{Validation selection and the sealed test}
Development headline accuracy values use Flickr30k validation. Multiple model choices touched that split, and the small \MTone{} gain in particular is vulnerable to selection pressure. The larger FreezeShift improvement is harder to explain away on magnitude alone, but it remains development evidence on that split.

The confirmatory one-shot result is reported in Section 4.13. It reduces but does not remove the threat from validation selection: only the two student endpoints were evaluated across seeds, while each reference remains a single fixed checkpoint without an accuracy-variance estimate. The sealed test confirms the model ordering and quantifies the strict-freezing cost, but it does not turn one dataset into evidence of broad semantic generalisation.

Three seeds quantify local run variability but provide weak tail uncertainty, and the deterministic reference evaluations do not estimate accuracy variance. The pooled-SD promotion rule is a decision heuristic, not a hypothesis test.

\section{Scale and generality}
The study uses a bounded set of encoders, COCO-centred training, one qualified CC3M mirror, and one principal hardware class. Conclusions about queues, data, and teachers are conditional on this setting. CC3M-only training did not converge before its compute ceiling, so no universal statement about web data is available. The final \MTone{} and FreezeShift endpoints are established for short-caption retrieval; Phase 2's mixed classification results do not establish their compositional, long-text, multilingual or domain-shift generalisation.

\section{Post-observation analyses}
The modality-specific drift regressions and latency repair were initiated after anomalies were observed. They are labelled accordingly. The queue regressions are descriptive and small-sample; the latency repair is stronger because its corrective protocol was frozen before a new allocation was measured.

\section{Mechanism identification}
The queue factorial separates several conditions but not false-negative structure, representation curvature, and gradient sensitivity. The dissertation establishes dissociation, not causation. Direct geometry-normalised drift and instrumented-gradient studies remain future work. FALCON~\cite{kim2026falcon} and PCME~\cite{chun2021pcme} make the false-negative alternative especially important: COCO and Flickr relevance is many-to-many. Known same-image repetitions were multi-positive masked; unlabelled cross-image semantic matches were not.

\section{Latency portability}
Absolute milliseconds depend on GPU, software, clocking, padding, batch, and input policy. The defensible quantity is the paired same-allocation comparison, not universal latency. TokenShift was accuracy-qualified on the declared validation protocol and rejected, but its exact trade-off remains specific to the tested fixed 2$\times$2 merge, depths, parents, and hardware.

\section{Literature coverage}
The novelty assessment is bounded to the screened corpus. The 10 August 2026 refresh adds recent direct or near-direct work including SOTAlign~\cite{roschmann2026sotalign}, HALoRA, FALCON, Dyna-ViT~\cite{rubab2026dynavit} and SAD-TM~\cite{xie2026sadtm}. Rapid work on frozen alignment and efficient probing still means the Tier-1 search should be repeated at formal submission. Absence from the corpus is not proof that no related result exists.

\section{Efficiency and governance coverage}
The repository measures parameters, FLOPs and latency, but it did not uniformly instrument peak training memory, energy, carbon or financial cost. Retrospective carbon estimates would introduce unrecorded assumptions and are not supplied. Likewise, the responsible-reporting section identifies inherited bias and licensing duties but is not a demographic fairness audit or legal opinion.

\chapter{Reproducibility, ethics, and responsible reporting}
The GitHub repository retains machine-readable experiment status, frozen configs, hashes, compact generated reports, and test coverage. Heavyweight datasets, checkpoints, caches, tensors, and Slurm logs are excluded from public submission but were moved into a locally recoverable, manifested quarantine. Failed and stopped studies remain represented by their configs and reports. A dependency lock records the verified environment specification.

The work uses established vision-language datasets and pretrained models. Its primary ethical risks arise from inherited dataset and model biases, especially in open-web captions and zero-shot classification. This dissertation does not audit demographic fairness and should not imply that retrieval performance establishes safe deployment. Dataset licences and pretrained-model licences must be checked for any redistribution or commercial use. The public submission does not redistribute Flickr30k/COCO images or third-party model weights; it retains identifiers, manifests, configuration and derived aggregate results. Software licences do not automatically grant rights in training images, captions, model weights or downstream uses.

Responsible reporting requires the following distinctions:

\begin{itemize}
\item validation versus test;
\item frozen backbones versus LoRA-adapted backbones;
\item trainable parameters versus total model parameters;
\item same-allocation relative latency versus portable absolute latency;
\item parameter-free token reduction versus a genuine accuracy--latency frontier improvement; and
\item an empirical boundary condition versus a proven mechanism.
\end{itemize}
\chapter{Conclusion}
This dissertation set out to determine whether independently pretrained vision and text encoders could approach a compact jointly pretrained VLM under tight adaptation and inference constraints. The answer is qualified: the work approaches the split-matched OpenCLIP field anchor under its local budget, but it does not demonstrate parity with MobileCLIP2-S0, the primary compact reference. That distinction is a result, not a relabelling opportunity. The answer emerged through correction as much as invention. Weak early results were not an immutable compatibility ceiling; a stale memory queue had caused most of the collapse. Historical embeddings were harmful immediately, and their damage could not be certified by drift magnitude alone. In-distribution coverage, recipe, padding, teacher choice, and measured rather than nominal efficiency all changed conclusions.

Within the fully frozen-backbone regime, learned token aggregation produced \MTone{} at 54.487\% validation R@1 with 2.896 million trainable parameters and faster same-allocation latency than OpenCLIP. Bounded LoRA adaptation raised the score to 63.416\% with 4.862 million trainable parameters, still faster than the reference under the paired measurement. This closes much of the gap but not all of it. The outcome is therefore neither parity nor failure: it is a quantified frontier, with an explicit cost for strict freezing and a measured benefit from limited adaptation.

The most durable contribution, though, is probably not the frontier. It is the queue result, and it is not confined to frozen alignment. Cross-batch memory is used well beyond this setting, and it is routinely justified by the argument that features move slowly enough for stored embeddings to remain usable. That argument turns out to be insufficient. Harm here began after a single optimizer step; conditions matched on measured drift differed in damage by up to 4.74 percentage points; and the text tower degraded roughly 7.2 times more steeply than the image tower over the same drift range. Drift and damage are associated, but drift does not determine damage, so a scalar drift threshold cannot certify that a stored negative is safe. The frozen encoders did not cause this. They made it measurable, by holding everything else still.

That result was only reachable because the programme was willing to be wrong in public. Queues, data scale, aggregation, adaptation and timing were each allowed to produce negative or inconvenient outcomes: a mining branch stopped by its own gate, a web-scale data source that failed to converge inside its compute ceiling, a latency optimisation that worked and was rejected because it cost too much accuracy, and a timing comparison that had to be declared invalid after it had already produced a verdict. The authorised one-shot final-test evaluation has since been carried out and closes the confirmatory sequence: on the sealed Flickr30k test the strictly frozen endpoint reaches 52.90 $\pm$ 0.87\% and dual FreezeShift 62.20 $\pm$ 0.45\% mean bidirectional R@1. Multi-seed reference evaluation, broader task generalisation and mechanism instrumentation remain legitimate future work rather than evidence silently assumed complete.

The honest summary is a narrow one. This work does not beat CLIP, and it does not reach compact-reference parity. What it establishes is how far efficient alignment of independently pretrained encoders can currently be pushed, what strict freezing costs when that cost is actually measured rather than assumed, and that a training component in common use can fail immediately, in modality-dependent ways, for reasons its standard justification does not predict.

\appendix
\chapter{Experiment inventory}
The authoritative inventory is \texttt{.\linebreak[1].\linebreak[1]/\linebreak[1]experiments/\linebreak[1]README.\linebreak[1]md}. The repository records 29 experiment families spanning foundations, adaptive reranking, recipe diagnosis, queue mechanism, data and transfer, efficiency and token architecture, diagnostics, bounded adaptation, and measurement repair. Phase 2, FreezeShift, TokenShift profiling, the latency amendment, and the TokenShift accuracy closure are complete. Alignment v3 remains incomplete-archived, the recovery branch is superseded, and guarded hard-negative mining is a valid stopped-by-gate result.

\chapter{Final evidence paths}
\begin{itemize}
\item Fully frozen \MTone{}: \texttt{results/\linebreak[1]text\_\linebreak[1]aggregation\_\linebreak[1]study/\linebreak[1]training/\linebreak[1]M\_\linebreak[1]T1/\linebreak[1]report/\linebreak[1]report.\linebreak[1]json}
\item Final LR decision: \texttt{results/\linebreak[1]final\_\linebreak[1]lr\_\linebreak[1]study/\linebreak[1]report/\linebreak[1]report.\linebreak[1]json}
\item FreezeShift: \texttt{freezeshift/\linebreak[1]results/\linebreak[1]report/\linebreak[1]report.\linebreak[1]json}
\item Corrective latency: \texttt{latency\_\linebreak[1]amendment/\linebreak[1]results/\linebreak[1]report/\linebreak[1]report.\linebreak[1]json}
\item Final model decision: \texttt{docs/\linebreak[1]final\_\linebreak[1]model\_\linebreak[1]freeze.\linebreak[1]md}
\item TokenShift: \texttt{tokenshift/\linebreak[1]results/\linebreak[1]profile/\linebreak[1]report.\linebreak[1]json}
\item TokenShift accuracy: \texttt{tokenshift\_\linebreak[1]training/\linebreak[1]results/\linebreak[1]report/\linebreak[1]report.\linebreak[1]json}
\item Queue factorial: \texttt{results/\linebreak[1]queue\_\linebreak[1]factorial/\linebreak[1]report/\linebreak[1]report.\linebreak[1]json}
\item Drift analysis: \texttt{results/\linebreak[1]queue\_\linebreak[1]factorial/\linebreak[1]interpretation/\linebreak[1]drift\_\linebreak[1]regression.\linebreak[1]json}
\item Known-positive audit: \texttt{results/\linebreak[1]queue\_\linebreak[1]factorial/\linebreak[1]interpretation/\linebreak[1]known\_\linebreak[1]positive\_\linebreak[1]audit.\linebreak[1]md}
\item Phase 2: \texttt{results/\linebreak[1]phase2/\linebreak[1]phase2\_\linebreak[1]report.\linebreak[1]json}
\item Experiment registry: \texttt{experiments/\linebreak[1]registry.\linebreak[1]yaml}
\end{itemize}
